\chapter{Fundamentarea teoretică și documentarea bibliografică}
\label{cap:cap2}

Tranzacționarea algoritmică poate fi definită ca procesul prin care deciziile de cumpărare, vânzare sau menținere a unei poziții sunt generate automat, pe baza unor reguli formale. Spre deosebire de tranzacționarea discreționară, în care decizia depinde în mare măsură de interpretarea subiectivă a operatorului, un sistem algoritmic aplică aceeași logică asupra fiecărei observații disponibile. Această proprietate este importantă deoarece permite reproducerea rezultatelor și analiza obiectivă a comportamentului strategiei.

În literatura de specialitate se subliniază faptul că o strategie de tranzacționare nu trebuie evaluată doar prin exemple izolate de tranzacții profitabile, ci printr-un proces riguros de testare, măsurare a riscului și verificare a robusteții \cite{book:chan:algo_trading}. În acest sens, prezenta lucrare tratează backtesting-ul, managementul riscului și raportarea rezultatelor ca elemente centrale, nu ca funcționalități secundare.

\section{Piața Forex și datele OHLCV}
\label{cap:cap2:forex}

Piața Forex reprezintă piața globală de schimb valutar, în care instrumentele sunt exprimate sub forma unor perechi valutare, precum EURUSD, GBPUSD sau USDJPY. Prima monedă din pereche poartă denumirea de monedă de bază, iar a doua monedă este moneda cotată. Prețul perechii indică numărul de unități din moneda cotată necesare pentru a cumpăra o unitate din moneda de bază. Această structură influențează modul de calcul al profitului și pierderii, mai ales atunci când moneda cotată nu este USD.

Datele utilizate în proiect sunt de tip OHLCV, adică includ prețul de deschidere, maximul, minimul, prețul de închidere și volumul pentru un interval de timp. Sistemul suportă mai multe granularități, de la intervale intraday până la timeframe-uri zilnice sau săptămânale. În contextul backtesting-ului, ordinea cronologică a acestor observații este esențială. Un algoritm care ar utiliza informații din viitor ar produce rezultate nerealiste, motiv pentru care toate componentele proiectului calculează indicatorii folosind doar date disponibile până la momentul curent.

\section{Particularități operaționale ale pieței Forex}
\label{cap:cap2:particularitati_forex}

Piața Forex are particularități care influențează proiectarea unui sistem automatizat. Spre deosebire de acțiuni, unde fiecare instrument este exprimat într-o singură monedă de referință, perechile valutare implică simultan o monedă de bază și o monedă cotată. Din acest motiv, calculul profitului, al pierderii și al expunerii trebuie să țină cont de conversia către moneda contului. În lucrare, raportarea este realizată în USD, iar pentru perechile unde USD nu este moneda cotată este necesară o rată de conversie suplimentară.

Un alt aspect important este faptul că piața funcționează aproape continuu în timpul săptămânii, dar lichiditatea nu este uniformă pe toate intervalele orare. Sesiunile asiatice, europene și americane pot avea comportamente diferite, atât prin volatilitate, cât și prin frecvența mișcărilor direcționale. O strategie care produce semnale bune într-un interval cu lichiditate ridicată poate deveni instabilă în perioade cu volum redus sau spread mai mare. Din acest motiv, sistemul include filtre de sesiune și verificări legate de volatilitate.

În practică, rezultatul unei tranzacții nu depinde doar de direcția prețului, ci și de costurile implicite ale execuției: spread, slippage și eventuale comisioane. Chiar dacă proiectul are caracter educațional și folosește ipoteze simplificate în backtesting, arhitectura păstrează posibilitatea configurării slippage-ului și a costurilor. Această separare este importantă deoarece permite extinderea ulterioară către simulări mai apropiate de condițiile reale de piață.

\section{Calitatea datelor și evitarea erorilor de simulare}
\label{cap:cap2:calitatea_datelor}

Calitatea datelor istorice influențează direct validitatea rezultatelor obținute prin backtesting. Datele lipsă, barele duplicate, ordonarea incorectă sau valorile extreme pot modifica indicatorii tehnici și pot genera semnale care nu ar fi apărut într-un flux real de piață. De aceea, un sistem de testare nu trebuie să trateze datele doar ca intrare brută, ci ca resursă care trebuie verificată înainte de utilizare.

În cazul seriilor OHLCV, validarea minimă presupune verificarea ordinii cronologice, a consistenței dintre valorile open, high, low și close, precum și eliminarea observațiilor imposibile. De exemplu, maximul unei bare trebuie să fie mai mare sau egal cu prețul de deschidere și închidere, iar minimul trebuie să fie mai mic sau egal cu acestea. Abaterile de la aceste reguli pot indica date corupte sau probleme de descărcare.

O altă sursă de eroare este folosirea accidentală a informației viitoare. Aceasta poate apărea atunci când indicatorii sunt calculați pe întreaga serie și apoi utilizați fără decalare, sau când deciziile curente includ valori care în realitate ar fi fost disponibile doar după închiderea barei. În proiect, procesarea bară cu bară și folosirea unui mecanism de replay cronologic reduc acest risc, deoarece fiecare decizie este generată pe baza istoricului disponibil până la momentul curent.

\section{Indicatori tehnici utilizați}
\label{cap:cap2:indicatori}

Strategiile implementate se bazează pe indicatori tehnici clasici, aleși pentru interpretabilitate și pentru compatibilitatea lor cu datele OHLCV. Scopul acestor indicatori nu este de a prezice perfect evoluția prețului, ci de a descrie anumite caracteristici ale pieței, precum direcția, volatilitatea sau poziționarea prețului față de o medie.

\textit{Media mobilă exponențială} (EMA) acordă o importanță mai mare observațiilor recente și este utilizată pentru identificarea direcției trendului. În strategia de trend, o EMA rapidă este comparată cu o EMA lentă, iar semnalul este confirmat suplimentar prin MACD și ADX.

\textit{Relative Strength Index} (RSI) măsoară raportul dintre mișcările recente ascendente și descendente. În cadrul strategiei de revenire la medie, valorile foarte mici ale RSI pot indica o zonă de supravânzare, iar valorile foarte mari pot indica supracumpărare.

\textit{Benzile Bollinger} definesc un canal statistic în jurul unei medii mobile. Apropierea prețului de banda inferioară sau de banda superioară poate semnala o abatere temporară față de nivelul mediu, aspect util în regimuri de piață laterale.

\textit{Canalele Donchian} utilizează maximele și minimele ultimelor bare pentru identificarea străpungerilor. Într-un regim de trend sau de volatilitate ridicată, depășirea canalului poate indica o posibilă continuare a mișcării.

\textit{Average True Range} (ATR) estimează volatilitatea și este folosit atât pentru validarea semnalelor, cât și pentru dimensionarea stop-loss-ului și a poziției. Prin raportarea riscului la ATR, sistemul adaptează dimensiunea tranzacției la condițiile curente de piață.

\textit{Average Directional Index} (ADX) este utilizat pentru evaluarea forței trendului, iar exponentul Hurst oferă o perspectivă suplimentară asupra caracterului persistent sau oscilant al seriei de preț.

\section{Regimuri de piață}
\label{cap:cap2:regimuri}

Un sistem de tranzacționare care aplică aceeași logică în toate condițiile de piață poate avea performanțe instabile. Piața poate evolua direcțional, lateral sau poate traversa perioade de volatilitate ridicată. Din acest motiv, proiectul clasifică fiecare bară într-un regim de piață: \textit{trend\_up}, \textit{trend\_down}, \textit{range} sau \textit{volatility\_spike}. Clasificarea se bazează pe indicatori precum ADX, ATR, z-score-ul volatilității, exponentul Hurst și poziția prețului față de o medie mobilă.

Această clasificare permite activarea strategiilor în condițiile pentru care au fost proiectate. Strategia de trend este relevantă în regimuri direcționale, strategia de revenire la medie este mai potrivită pentru piețe laterale, iar strategia de breakout devine utilă atunci când volatilitatea crește sau când prețul depășește intervale recente importante.

\section{Backtesting și riscul de supraoptimizare}
\label{cap:cap2:backtesting}

Backtesting-ul reprezintă simularea aplicării unei strategii pe date istorice. Acesta este un pas necesar în evaluarea inițială a unei strategii, dar nu constituie o garanție pentru performanța viitoare. O problemă frecventă este supraoptimizarea, adică ajustarea excesivă a parametrilor pentru un set istoric particular. Într-o astfel de situație, strategia poate obține rezultate bune pe perioada testată, dar poate eșua atunci când este aplicată pe date noi. Bailey et al. evidențiază necesitatea unei evaluări prudente, mai ales atunci când sunt testate multe variante ale aceleiași strategii \cite{art:bailey:backtest}.

În cadrul proiectului, backtesting-ul este construit determinist. Datele sunt încărcate din cache local, fiecare bară este procesată o singură dată, iar rezultatele sunt salvate într-un director dedicat. Această abordare permite reluarea experimentelor și verificarea tranzacțiilor generate, aspect important pentru o lucrare tehnică în care rezultatele trebuie să poată fi justificate.

\section{Metrici de evaluare a performanței}
\label{cap:cap2:metrici_performanta}

Evaluarea unei strategii de tranzacționare nu poate fi realizată doar prin profitul final. O strategie poate obține un profit ridicat, dar cu un drawdown foarte mare, ceea ce o face dificil de utilizat în practică. În mod similar, o rată mare a tranzacțiilor câștigătoare nu garantează profitabilitatea dacă pierderile medii sunt mai mari decât câștigurile medii. De aceea, analiza trebuie să combine indicatori de randament, risc și stabilitate.

\textit{Profit factor} este raportul dintre profitul brut și pierderea brută. O valoare mai mare decât 1 indică faptul că profiturile totale depășesc pierderile totale, însă interpretarea trebuie făcută împreună cu numărul de tranzacții. Un profit factor ridicat obținut din foarte puține tranzacții poate fi mai puțin convingător decât o valoare moderată obținută pe un eșantion mai mare.

\textit{Drawdown-ul maxim} arată scăderea maximă a capitalului față de un vârf anterior. Această metrică este importantă deoarece descrie presiunea pe care strategia o pune asupra capitalului și asupra disciplinei operaționale. Într-un sistem automatizat, drawdown-ul este folosit nu doar pentru evaluare, ci și pentru mecanisme de protecție, precum circuit breaker-ul.

\textit{Sharpe ratio} și \textit{Sortino ratio} sunt indicatori de performanță ajustată la risc. Sharpe raportează randamentul la volatilitatea totală, în timp ce Sortino penalizează mai ales volatilitatea negativă. În contextul tranzacționării, Sortino poate fi mai relevant deoarece fluctuațiile pozitive nu reprezintă un risc în același sens ca pierderile.

Comparația cu un reper simplu, precum buy-and-hold, este utilă pentru interpretarea rezultatelor. Dacă o strategie activă produce un profit asemănător cu menținerea pasivă a poziției, dar cu un risc mai mare sau cu o complexitate operațională mai ridicată, valoarea sa practică este discutabilă. Din acest motiv, lucrarea raportează rezultatele sistemului împreună cu o comparație față de comportamentul pasiv al perechii valutare.

\section{Managementul riscului}
\label{cap:cap2:risc}

Managementul riscului are un rol central în sistemele de tranzacționare, deoarece rezultatul unei strategii depinde nu doar de calitatea semnalelor, ci și de modul în care sunt controlate expunerea, volatilitatea și pierderile succesive. O strategie poate genera semnale corecte din punct de vedere tehnic, dar poate produce pierderi semnificative dacă dimensiunea pozițiilor nu este controlată. În proiect, fiecare intrare este validată de o componentă de risc independentă, care ține cont de capitalul disponibil, volatilitate, expunere totală și drawdown.

Regula implicită presupune riscul a 1\% din capital pentru o tranzacție, iar expunerea totală este limitată. Stop-loss-ul este calculat ca multiplu de ATR, iar take-profit-ul utilizează o distanță adaptată volatilității. Sistemul include, de asemenea, mecanisme de protecție precum break-even, trailing stop și circuit breaker. Dacă drawdown-ul depășește pragul configurat, tranzacționarea este dezactivată temporar, pentru a preveni acumularea rapidă a pierderilor.

\section{Soluții similare și poziționarea proiectului}
\label{cap:cap2:solutii}

Există numeroase soluții pentru tranzacționare algoritmică, de la platforme comerciale consacrate până la framework-uri open-source pentru cercetare cantitativă. Platformele integrate, precum MetaTrader, oferă conectare directă la broker și un ecosistem matur, însă pot introduce dependențe de platformă, de limbajul specific și de calitatea infrastructurii brokerului. Framework-urile Python, precum Backtrader sau Zipline, sunt flexibile și potrivite pentru cercetare, dar necesită adesea adaptări pentru integrare live, raportare și gestionarea particularităților pieței Forex. Platformele cloud, precum QuantConnect, reduc efortul operațional, dar introduc dependențe externe și costuri suplimentare.

\begin{table}[htbp]
    \centering
    \caption{Comparație între tipuri de soluții pentru tranzacționare algoritmică}
    \label{tabel:comparatie_solutii}
    \begin{tabular}{|p{3.3cm}|p{4.2cm}|p{4.2cm}|p{3.6cm}|}
        \hline
        \textbf{Soluție} & \textbf{Avantaje} & \textbf{Limitări} & \textbf{Poziționarea proiectului} \\
        \hline
        MetaTrader / roboți Expert Advisor & Integrare directă cu brokeri, interfață matură, ecosistem larg. & Dependență de platformă, limbaj specific, control mai redus asupra fluxului de cercetare. & Proiectul oferă cod Python transparent și ușor de auditat. \\
        \hline
        Framework-uri Python de backtesting & Flexibilitate, biblioteci numerice mature, potrivite pentru cercetare. & Necesită integrare separată cu brokeri și raportare personalizată. & Proiectul include atât backtesting, cât și integrare paper cu Interactive Brokers. \\
        \hline
        Platforme cloud de trading algoritmic & Infrastructură administrată, date și execuție integrate. & Costuri, dependență de furnizor și control mai redus asupra mediului local. & Proiectul rulează local, cu date cache-uite și rezultate reproductibile. \\
        \hline
        Scripturi simple de semnalizare & Ușor de implementat și modificat. & Lipsă de arhitectură, risc slab controlat, raportare incompletă. & Proiectul separă datele, strategiile, riscul, execuția și raportarea. \\
        \hline
    \end{tabular}
\end{table}

Soluția propusă se poziționează ca un proiect educațional și experimental, orientat către claritatea fluxului decizional. Sistemul nu urmărește execuție de frecvență înaltă și nu promite profitabilitate, ci oferă un cadru prin care strategiile pot fi evaluate, comparate și documentate. Valoarea principală a proiectului constă în faptul că fiecare etapă a deciziei poate fi urmărită: de la datele de intrare, până la semnal, validarea riscului și rezultatul final al tranzacției.

Lacuna pe care proiectul o acoperă este lipsa unei soluții compacte, ușor de inspectat, care să combine într-un singur cod sursă descărcarea și cache-ul datelor Forex, backtesting determinist, strategii complementare, allocator explicabil, management de risc independent, integrare paper trading și export de artefacte pentru analiză.

\section{Cerințe derivate}
\label{cap:cap2:cerinte}

Din analiza domeniului rezultă următoarele cerințe pentru aplicație:
\begin{itemize}
    \item datele istorice trebuie să fie curate, ordonate cronologic și reutilizabile;
    \item strategiile trebuie să fie independente și ușor de extins;
    \item decizia finală trebuie să fie explicabilă prin scorurile strategiilor participante;
    \item riscul trebuie aplicat înainte de deschiderea fiecărei poziții;
    \item rezultatele trebuie salvate în formate care permit auditarea tranzacțiilor;
    \item modul live trebuie să includă protecții împotriva rulării accidentale pe conturi reale;
    \item monitorizarea live trebuie să acopere prețurile, semnalele, pozițiile și performanța curentă.
\end{itemize}

Prin realizarea acestei analize comparative și a documentării bibliografice, a fost atins primul obiectiv al lucrării, respectiv studiul literaturii de specialitate și identificarea cerințelor relevante pentru un sistem de tranzacționare algoritmică. Analiza confirmă necesitatea dezvoltării unei soluții integrate, reproductibile și ușor de auditat, care să combine backtesting-ul determinist, controlul riscului, raportarea rezultatelor și verificarea controlată în regim paper trading.
